\chapter{Measure-Theoretic Foundations}
\label{ch:measure}

Up to now we have used probability spaces $(\Omega,\F,\Prob)$, random variables,
expectations and conditional expectations, and we have occasionally deferred a
technical point to this chapter. Here we collect and prove those points: why
$\sigma$-algebras are needed, what a random variable \emph{is}, how expectation is
defined as an integral, the three convergence theorems (monotone, Fatou, dominated),
Fubini's theorem, independence and Kolmogorov's $0$--$1$ law, the existence of independent
sequences, the characteristic function theorems (inversion, uniqueness, L\'evy continuity),
and the rigorous definition of conditional expectation.

The chapter assumes the real-analysis background of sequences, limits and suprema; where
a deep external theorem is needed (Carath\'eodory's extension theorem, Helly's selection
theorem, the Riesz representation theorem), it is stated and clearly marked as cited.

\section{$\sigma$-algebras and measurable functions}

Recall (Chapter~\ref{ch:axioms}) that a $\sigma$-algebra $\F$ on $\Omega$ contains
$\Omega$ and is closed under complements and countable unions. The intersection of any
collection of $\sigma$-algebras on $\Omega$ is again a $\sigma$-algebra. Hence for any
collection $\mathcal C$ of subsets there is a smallest $\sigma$-algebra containing it, the
\emph{$\sigma$-algebra generated by $\mathcal C$}, written $\sigma(\mathcal C)$ (the
intersection of all $\sigma$-algebras containing $\mathcal C$).

\begin{definition}[Borel $\sigma$-algebra]
The \emph{Borel $\sigma$-algebra} $\mathcal B(\R)$ is generated by the open sets of $\R$.
Equivalently it is generated by the intervals $(-\infty,x]$, $x\in\R$ (every open set is a
countable union of open intervals, and each open interval $(a,b)$ is obtained from
half-lines by countable unions, intersections and complements).
\end{definition}

\begin{definition}[Measurable function; random variable]
Let $(\Omega,\F)$ and $(E,\mathcal E)$ be measurable spaces. A function $f:\Omega\to E$ is
\emph{measurable} if $f^{-1}(B)\in\F$ for every $B\in\mathcal E$. A \emph{random variable} is
a measurable function $X:(\Omega,\F)\to(\R,\mathcal B(\R))$. A \emph{random vector} is a
measurable map into $(\R^d,\mathcal B(\R^d))$.
\end{definition}

\begin{proposition}\label{prop:measurable}
\begin{enumerate}[label=(\roman*)]
\item $f$ is measurable iff $f^{-1}(C)\in\F$ for all $C$ in a collection $\mathcal C$ that
generates $\mathcal E$. In particular $X:\Omega\to\R$ is a random variable iff
$\{X\le x\}\in\F$ for every $x$ (the definition used in Chapter~\ref{ch:drv}).
\item Compositions of measurable functions are measurable; continuous functions
$\R\to\R$ are Borel measurable.
\item If $X,Y$ are random variables and $c\in\R$, then $X+Y$, $XY$, $cX$, $\max(X,Y)$ are
random variables; if $X_n$ are random variables then $\sup_nX_n$, $\inf_nX_n$,
$\limsup X_n$, $\liminf X_n$ are random variables (possibly $\pm\infty$), and
$\{X_n\text{ converges}\}\in\F$.
\end{enumerate}
\end{proposition}
\begin{proof}
(i) The collection $\{B\in\mathcal E:f^{-1}(B)\in\F\}$ is a $\sigma$-algebra (preimages
commute with complements and unions) containing $\mathcal C$, hence contains
$\sigma(\mathcal C)=\mathcal E$. (ii) Preimages of preimages; the preimage of an open set under
a continuous function is open. (iii) $\{X+Y<c\}=\bigcup_{q\in\mathbb Q}\{X<q\}\cap\{Y<c-q\}$ is a
countable union of events; the other cases are similar. $\{\sup X_n\le x\}=\bigcap_n\{X_n\le x\}$,
$\liminf X_n=\sup_n\inf_{m\ge n}X_m$, and $\{X_n\text{ converges}\}=\{\liminf X_n=\limsup X_n\}$.
\end{proof}

\begin{definition}[Generated $\sigma$-algebra, law]
For a random variable $X$, $\sigma(X)=\{X^{-1}(B):B\in\mathcal B(\R)\}$ is the smallest
$\sigma$-algebra making $X$ measurable; it represents ``the information in $X$''. The
\emph{law} (or distribution) of $X$ is the probability measure
$\mu_X(B)=\Prob(X\in B)$ on $(\R,\mathcal B(\R))$.
\end{definition}

\begin{nuance}{measurability is a statement about information}
A function $Z$ is $\sigma(X)$-measurable iff $Z=g(X)$ for some Borel function $g$ (the
Doob--Dynkin lemma, cited). So a random variable is determined by $X$ exactly when it is
a function of $X$. This is the precise meaning of ``$\F_n$ is the information in
$X_1,\dots,X_n$'' (Chapter~\ref{ch:mart}) and of ``determined by $X_1,\dots,X_n$''.
\end{nuance}

\section{Measures, the $\pi$--$\lambda$ theorem, and the Vitali set}

\begin{definition}[Measure]
A \emph{measure} on $(\Omega,\F)$ is a function $\mu:\F\to[0,\infty]$ with
$\mu(\emptyset)=0$ and countable additivity for disjoint sets. It is \emph{finite} if
$\mu(\Omega)<\infty$, a \emph{probability} measure if $\mu(\Omega)=1$, and
\emph{$\sigma$-finite} if $\Omega=\bigcup_n\Omega_n$ with $\mu(\Omega_n)<\infty$.
\end{definition}

The properties of Chapter~\ref{ch:axioms} (monotonicity, countable subadditivity,
continuity from below, and from above when $\mu$ is finite) hold for measures with the same
proofs. The basic example beyond probability is \emph{Lebesgue measure} $\lambda$ on
$\R$, the unique measure on $\mathcal B(\R)$ with $\lambda((a,b])=b-a$. Its existence is
Carath\'eodory's extension theorem (cited): a countably additive function on an algebra of
sets extends to the generated $\sigma$-algebra, and uniquely if $\sigma$-finite.

\begin{theorem}[Existence of non-measurable sets]\label{thm:vitali}
There is no countably additive, translation-invariant measure $\lambda$ on \emph{all}
subsets of $[0,1)$ with $\lambda([0,1))=1$. Hence a $\sigma$-algebra strictly smaller than
$2^{[0,1)}$ is unavoidable.
\end{theorem}
\begin{proof}
Suppose $\lambda$ exists, with translation understood modulo $1$ (so that translation
preserves $[0,1)$). Define $x\sim y$ if $x-y\in\mathbb Q$; this is an equivalence relation on
$[0,1)$. By the axiom of choice there is a set $V$ containing exactly one point of each
equivalence class. For each rational $q\in[0,1)$ let $V_q=V+q\pmod1$. The sets $V_q$ are
pairwise disjoint (a common point would put two points of $V$ in one class) and their union
is $[0,1)$ (every $x$ is equivalent to some $v\in V$, so $x=v+q$ for some rational
$q$ modulo $1$). Translation invariance gives $\lambda(V_q)=\lambda(V)=c$ for all $q$, so by
countable additivity $1=\sum_{q}c$, which is $0$ if $c=0$ and $\infty$ if $c>0$. Contradiction.
\end{proof}

\begin{definition}[$\pi$-system and $\lambda$-system]
A collection $\mathcal P$ of subsets of $\Omega$ is a \emph{$\pi$-system} if it is closed
under finite intersections. A collection $\mathcal L$ is a \emph{$\lambda$-system} if
(i) $\Omega\in\mathcal L$, (ii) $A,B\in\mathcal L$ and $A\subseteq B$ imply $B\setminus A\in\mathcal L$,
and (iii) increasing unions of members of $\mathcal L$ are in $\mathcal L$.
\end{definition}

\begin{theorem}[Dynkin's $\pi$--$\lambda$ theorem]\label{thm:pilambda}
If $\mathcal P$ is a $\pi$-system, $\mathcal L$ a $\lambda$-system, and $\mathcal P\subseteq\mathcal L$,
then $\sigma(\mathcal P)\subseteq\mathcal L$.
\end{theorem}
\begin{proof}
Let $\ell(\mathcal P)$ be the smallest $\lambda$-system containing $\mathcal P$
(the intersection of all of them). It suffices to show $\ell(\mathcal P)$ is a
$\sigma$-algebra, since then $\sigma(\mathcal P)\subseteq\ell(\mathcal P)\subseteq\mathcal L$.
A $\lambda$-system is closed under complements ($A^c=\Omega\setminus A$), so it is a
$\sigma$-algebra iff it is closed under finite intersections (then finite unions follow by
De Morgan, and increasing limits give countable unions).

For $A\in\mathcal P$ let $\mathcal L_A=\{B:A\cap B\in\ell(\mathcal P)\}$. It is a $\lambda$-system
($A\cap\Omega=A$; if $B_1\subseteq B_2$ in $\mathcal L_A$ then
$A\cap(B_2\setminus B_1)=(A\cap B_2)\setminus(A\cap B_1)$; increasing unions are preserved) and
contains $\mathcal P$ (as $\mathcal P$ is a $\pi$-system). So $\mathcal L_A\supseteq\ell(\mathcal P)$:
$A\cap B\in\ell(\mathcal P)$ for $A\in\mathcal P$, $B\in\ell(\mathcal P)$. Now for $B\in\ell(\mathcal P)$
let $\mathcal L_B'=\{A:A\cap B\in\ell(\mathcal P)\}$, again a $\lambda$-system; by the first step it
contains $\mathcal P$, hence contains $\ell(\mathcal P)$. Thus $\ell(\mathcal P)$ is closed under
intersections.
\end{proof}

\begin{corollary}[Uniqueness of measures]\label{cor:unique}
If two probability measures agree on a $\pi$-system $\mathcal P$, they agree on
$\sigma(\mathcal P)$. In particular, a probability measure on $\mathcal B(\R)$ is determined
by its CDF (the half-lines $(-\infty,x]$ form a $\pi$-system generating $\mathcal B(\R)$); and
the law of a random vector $(X_1,\dots,X_n)$ is determined by the joint CDF.
\end{corollary}
\begin{proof}
The collection of events on which the measures agree contains $\Omega$, is closed under
proper differences (subtract probabilities) and increasing unions (continuity from below),
so it is a $\lambda$-system containing $\mathcal P$. Apply Theorem~\ref{thm:pilambda}.
\end{proof}

This corollary justifies the claim, used throughout, that the CDF determines the
distribution.

\section{The Lebesgue integral and expectation}

For a measure space $(\Omega,\F,\mu)$ and $f\ge0$ measurable, the integral is built in three
stages.

\begin{enumerate}[label=\textbf{Stage \arabic*.}]
\item A \emph{simple function} is $s=\sum_{i=1}^ma_i\indic_{A_i}$ with $a_i\ge0$ and
$A_i\in\F$ disjoint; define $\int s\,d\mu=\sum_ia_i\mu(A_i)$ (independent of the
representation).
\item For measurable $f\ge0$, $\int f\,d\mu=\sup\bigl\{\int s\,d\mu:\ s\text{ simple},\,s\le f\bigr\}\in[0,\infty]$.
\item For general $f$ write $f=f^+-f^-$ with $f^\pm=\max(\pm f,0)$. If $\int f^+<\infty$ and
$\int f^-<\infty$ (equivalently $\int\abs f<\infty$), $f$ is \emph{integrable} and
$\int f\,d\mu=\int f^+d\mu-\int f^-d\mu$.
\end{enumerate}

For a random variable $X$ on $(\Omega,\F,\Prob)$, $\Exp[X]:=\int X\,d\Prob$. For the discrete and
density cases of earlier chapters this agrees with the sums and integrals defined there
(Theorem~\ref{thm:lotusgen} below). Monotonicity ($f\le g\Rightarrow\int f\le\int g$) and
$\int cf=c\int f$ for $c\ge0$ follow directly from Stage~2.

\begin{theorem}[Monotone convergence theorem, MCT]\label{thm:mct}
If $0\le f_1\le f_2\le\dots$ are measurable and $f_n\uparrow f$ pointwise, then
$\int f_n\,d\mu\uparrow\int f\,d\mu$.
\end{theorem}
\begin{proof}
Monotonicity gives $\int f_n\le\int f_{n+1}\le\int f$, so $L=\lim\int f_n\le\int f$. For the
reverse, let $s\le f$ be simple and $0<c<1$. The sets $E_n=\{f_n\ge cs\}$ increase to $\Omega$
(if $s(\omega)=0$ trivially; if $s(\omega)>0$ then $cs(\omega)<s(\omega)\le f(\omega)$ and
$f_n(\omega)\uparrow f(\omega)$). Then $\int f_n\ge\int f_n\indic_{E_n}\ge c\int s\indic_{E_n}
=c\sum_ia_i\mu(A_i\cap E_n)\to c\sum_ia_i\mu(A_i)=c\int s$ by continuity of measure from
below. So $L\ge c\int s$; let $c\uparrow1$ and take the supremum over $s$.
\end{proof}

\begin{theorem}[Linearity]\label{thm:linearity}
If $f,g$ are integrable and $a,b\in\R$ then $af+bg$ is integrable and
$\int(af+bg)=a\int f+b\int g$. In particular, for random variables,
$\Exp[aX+bY]=a\Exp[X]+b\Exp[Y]$ whenever $X,Y$ are integrable.
\end{theorem}
\begin{proof}
For simple functions additivity is clear. For measurable $f,g\ge0$, choose simple
$s_n\uparrow f$, $t_n\uparrow g$ (e.g.\ $s_n=\min(n,2^{-n}\lfloor2^nf\rfloor)$). Then
$s_n+t_n\uparrow f+g$ and by the MCT
$\int(f+g)=\lim\int(s_n+t_n)=\lim\bigl(\int s_n+\int t_n\bigr)=\int f+\int g$. For
integrable $f,g$ of any sign, $(f+g)^++f^-+g^-=(f+g)^-+f^++g^+$; integrate both sides
using additivity for non-negative functions and rearrange (all terms are finite). Scalar
multiples are handled by $\int cf=c\int f$ for $c\ge0$ and $\int(-f)=-\int f$.
\end{proof}

\begin{theorem}[Fatou's lemma]\label{thm:fatou}
If $f_n\ge0$ are measurable then $\int\liminf f_n\,d\mu\le\liminf\int f_n\,d\mu$.
\end{theorem}
\begin{proof}
Let $g_n=\inf_{k\ge n}f_k$. Then $0\le g_n\uparrow\liminf f_n$ and $g_n\le f_k$ for $k\ge n$,
so $\int g_n\le\inf_{k\ge n}\int f_k$. By MCT, $\int\liminf f_n=\lim\int g_n\le\liminf\int f_n$.
\end{proof}

\begin{theorem}[Dominated convergence theorem, DCT]\label{thm:dct}
If $f_n\to f$ pointwise (or almost everywhere), and $\abs{f_n}\le g$ for an integrable $g$,
then $f$ is integrable and $\int f_n\to\int f$; in fact $\int\abs{f_n-f}\to0$.
\end{theorem}
\begin{proof}
Apply Fatou to $g-f_n\ge0$ and $g+f_n\ge0$:
$\int(g-f)\le\liminf\int(g-f_n)=\int g-\limsup\int f_n$ and
$\int(g+f)\le\int g+\liminf\int f_n$. Subtracting the (finite) $\int g$ gives
$\limsup\int f_n\le\int f\le\liminf\int f_n$. For the last claim apply the first to
$\abs{f_n-f}\le2g$.
\end{proof}

\begin{theorem}[Change of variables (LOTUS)]\label{thm:lotusgen}
If $X$ has law $\mu_X$ and $g:\R\to\R$ is Borel with $g\ge0$ or $g(X)$ integrable, then
$\Exp[g(X)]=\int_\R g\,d\mu_X$. If $\mu_X$ has density $f$ with respect to Lebesgue measure
then $\Exp[g(X)]=\int g(x)f(x)\,dx$, and if $X$ is discrete, $\Exp[g(X)]=\sum_xg(x)p_X(x)$.
\end{theorem}
\begin{proof}
For $g=\indic_B$: $\Exp[\indic_B(X)]=\Prob(X\in B)=\mu_X(B)=\int\indic_Bd\mu_X$. By linearity
the formula holds for simple $g$, by MCT for $g\ge0$ (take simple $g_n\uparrow g$), and by
splitting $g=g^+-g^-$ in general. The density and discrete forms are the special cases
$d\mu_X=f\,dx$ and $\mu_X=\sum p_X(x)\delta_x$.
\end{proof}

\begin{definition}[Almost sure properties and $L^p$]
A property holds \emph{almost surely} (a.s.) if the exceptional set has probability zero.
If $f\ge0$ and $\int f=0$ then $f=0$ a.s.\ (since $\mu(f>1/n)\le n\int f=0$). Changing a
function on a null set does not change its integral. $L^p$ denotes the space of random
variables with $\Exp\abs X^p<\infty$ (modulo almost sure equality).
\end{definition}

\begin{theorem}[Tonelli and Fubini]\label{thm:fubini}
Let $(\Omega_1,\F_1,\mu_1)$, $(\Omega_2,\F_2,\mu_2)$ be $\sigma$-finite measure spaces and
$\mu_1\otimes\mu_2$ the product measure on $\F_1\otimes\F_2=\sigma(A\times B)$, satisfying
$(\mu_1\otimes\mu_2)(A\times B)=\mu_1(A)\mu_2(B)$. For a measurable $f:\Omega_1\times\Omega_2\to\R$:
\begin{enumerate}[label=(\roman*)]
\item (\emph{Tonelli}) if $f\ge0$ then $\int\!\!\int f\,d\mu_2\,d\mu_1=\int\!\!\int f\,d\mu_1\,d\mu_2=\int f\,d(\mu_1\otimes\mu_2)$;
\item (\emph{Fubini}) if $\int\abs f\,d(\mu_1\otimes\mu_2)<\infty$ the same equalities hold.
\end{enumerate}
\end{theorem}
\begin{proof}[Proof (finite measures)]
Let $\mathcal D$ be the class of sets $C\in\F_1\otimes\F_2$ for which $x\mapsto\mu_2(C_x)$ is
measurable (where $C_x=\{y:(x,y)\in C\}$) and $\int\mu_2(C_x)d\mu_1(x)=(\mu_1\otimes\mu_2)(C)$.
It contains the rectangles $A\times B$, a $\pi$-system generating $\F_1\otimes\F_2$, and it is
a $\lambda$-system: closed under proper differences and increasing unions (by linearity and
the MCT). By Theorem~\ref{thm:pilambda}, $\mathcal D=\F_1\otimes\F_2$, so (i) holds for
indicators, hence for simple functions by linearity, and for $f\ge0$ by the MCT. (The
$\sigma$-finite case follows by exhausting $\Omega_i$ with finite-measure sets, and the
other order of integration is symmetric.) For (ii) apply (i) to $f^+$ and $f^-$.
\end{proof}

\section{Independence and the zero--one law}

\begin{definition}[Independence of $\sigma$-algebras and variables]
Sub-$\sigma$-algebras $\mathcal G_1,\mathcal G_2,\dots\subseteq\F$ are \emph{independent} if
$\Prob(G_{i_1}\cap\dots\cap G_{i_k})=\prod_j\Prob(G_{i_j})$ for all distinct indices and
$G_i\in\mathcal G_i$. Random variables are independent iff the $\sigma$-algebras
$\sigma(X_i)$ are. Events are independent iff their indicator $\sigma$-algebras are.
\end{definition}

\begin{proposition}\label{prop:indpi}
If $\mathcal P_1,\mathcal P_2$ are $\pi$-systems with $\Prob(A\cap B)=\Prob(A)\Prob(B)$ for all
$A\in\mathcal P_1$, $B\in\mathcal P_2$, then $\sigma(\mathcal P_1)$ and $\sigma(\mathcal P_2)$ are
independent. Consequently random variables $X,Y$ are independent iff
$\Prob(X\le x,Y\le y)=\Prob(X\le x)\Prob(Y\le y)$ for all $x,y$, and iff the joint law is the
product law $\mu_X\otimes\mu_Y$.
\end{proposition}
\begin{proof}
Fix $A\in\mathcal P_1$. The class of $B$ with $\Prob(A\cap B)=\Prob(A)\Prob(B)$ is a
$\lambda$-system containing $\mathcal P_2$, hence contains $\sigma(\mathcal P_2)$. Now fix
$B\in\sigma(\mathcal P_2)$ and repeat for the class of $A$. The statement about CDFs takes
the $\pi$-systems of half-lines. The joint law and $\mu_X\otimes\mu_Y$ agree on the
rectangles $(-\infty,x]\times(-\infty,y]$ and hence everywhere by
Corollary~\ref{cor:unique}.
\end{proof}

\begin{theorem}[Product of independent variables]\label{thm:indprod}
If $X,Y$ are independent and integrable (or non-negative), then $\Exp[XY]=\Exp[X]\Exp[Y]$,
and $f(X),g(Y)$ are independent for Borel $f,g$.
\end{theorem}
\begin{proof}
The independence of $f(X)$ and $g(Y)$ holds since $\sigma(f(X))\subseteq\sigma(X)$ and
$\sigma(g(Y))\subseteq\sigma(Y)$. Let $X,Y\ge0$. The law of $(X,Y)$ is $\mu_X\otimes\mu_Y$, so
by Theorem~\ref{thm:lotusgen} and Tonelli,
$\Exp[XY]=\iint xy\,d\mu_X(x)d\mu_Y(y)=\int x\,d\mu_X\int y\,d\mu_Y$. For general integrable $X,Y$
split into positive and negative parts and use Fubini and linearity.
\end{proof}

\begin{theorem}[Kolmogorov's zero--one law]\label{thm:kolmogorov01}
Let $X_1,X_2,\dots$ be independent random variables and let
$\mathcal T=\bigcap_{n\ge1}\sigma(X_n,X_{n+1},\dots)$ be the \emph{tail $\sigma$-algebra}. Then
every $A\in\mathcal T$ has $\Prob(A)\in\{0,1\}$.
\end{theorem}
\begin{proof}
Fix $n$. By independence, $\sigma(X_1,\dots,X_n)$ is independent of
$\sigma(X_{n+1},X_{n+2},\dots)\supseteq\mathcal T$ (the independence of $\sigma$-algebras
generated by disjoint groups of independent variables follows from
Proposition~\ref{prop:indpi} applied to $\pi$-systems of finite-dimensional events). So $A\in\mathcal T$ is independent of
$\bigcup_n\sigma(X_1,\dots,X_n)$, a $\pi$-system that generates $\sigma(X_1,X_2,\dots)$. By
Proposition~\ref{prop:indpi}, $A$ is independent of $\sigma(X_1,X_2,\dots)$. But $A$ itself lies
in this $\sigma$-algebra, so $\Prob(A)=\Prob(A\cap A)=\Prob(A)^2$, whence $\Prob(A)\in\{0,1\}$.
\end{proof}

\begin{example}[Applications]
Events such as $\{\sum X_n\text{ converges}\}$, $\{X_n\text{ converges}\}$,
$\{\limsup S_n/n>c\}$ and $\{A_n\text{ i.o.}\}$ (for independent events $A_n$) are tail events:
changing finitely many $X_n$ does not affect them. So each has probability $0$ or $1$. This
explains the all-or-nothing behaviour in the Borel--Cantelli lemmas, and shows that a
random series either converges almost surely or diverges almost surely. The law
does not tell which: that is the role of Kolmogorov's three-series theorem (cited).
\end{example}

\subsection*{Existence of independent sequences}

All results about i.i.d.\ sequences presuppose that such sequences exist on some probability
space. Here is a self-contained construction.

\begin{theorem}[An i.i.d.\ sequence from one uniform variable]\label{thm:iidexist}
On $\Omega=[0,1)$ with Lebesgue measure there exist independent random variables
$U_1,U_2,\dots$, each uniform on $[0,1)$. Consequently for any distribution functions
$F_1,F_2,\dots$ there are independent $X_j=F_j^{-1}(U_j)$ with CDF $F_j$.
\end{theorem}
\begin{proof}
For $\omega\in[0,1)$ write its binary expansion $\omega=\sum_{k\ge1}\epsilon_k(\omega)2^{-k}$
(non-terminating expansions excluded). Each $\epsilon_k$ is measurable, and for any pattern
$(e_1,\dots,e_n)\in\{0,1\}^n$ the set $\{\epsilon_1=e_1,\dots,\epsilon_n=e_n\}$ is an interval
of length $2^{-n}$. Hence the $\epsilon_k$ are i.i.d.\ Bernoulli$(\tfrac12)$. Choose a bijection
$\sigma:\N\times\N\to\N$ and set $U_j=\sum_{i\ge1}\epsilon_{\sigma(i,j)}2^{-i}$. The $U_j$ are
functions of disjoint sub-families of independent bits, hence independent (by
Proposition~\ref{prop:indpi} on cylinder events), and each $U_j$ has the same law as
$\sum_i\epsilon_i2^{-i}$, which is uniform because
$\Prob(U_j\in[m2^{-n},(m+1)2^{-n}))=2^{-n}$ for all dyadic intervals, and dyadic intervals
form a $\pi$-system generating the Borel sets (Corollary~\ref{cor:unique}). The last
claim is Theorem~\ref{thm:pit} (and the quantile remark).
\end{proof}

(For arbitrary joint laws on infinite product spaces one uses Kolmogorov's extension
theorem, cited.)

\section{Absolutely continuous laws, Radon--Nikodym and the structure of distributions}

\begin{definition}[Absolute continuity]
$\nu$ is \emph{absolutely continuous} with respect to $\mu$, written $\nu\ll\mu$, if
$\mu(A)=0$ implies $\nu(A)=0$.
\end{definition}

\begin{theorem}[Radon--Nikodym]\label{thm:rn}
Let $\mu,\nu$ be finite measures on $(\Omega,\F)$ with $\nu\ll\mu$. Then there is a measurable
$h\ge0$, unique $\mu$-a.e., with $\nu(A)=\int_Ah\,d\mu$ for all $A\in\F$. We write $h=d\nu/d\mu$.
\end{theorem}
\begin{proof}[Proof (von Neumann)]
Let $\rho=\mu+\nu$. The map $\ell(f)=\int f\,d\nu$ is a bounded linear functional on
$L^2(\rho)$ since $\abs{\int f\,d\nu}\le\int\abs f\,d\rho\le\sqrt{\rho(\Omega)}\norm f_{L^2(\rho)}$.
By the Riesz representation theorem for Hilbert spaces (cited) there is $g\in L^2(\rho)$ with
$\int f\,d\nu=\int fg\,d\rho$ for all $f\in L^2(\rho)$. Testing with $f=\indic_{\{g<0\}}$ and
$f=\indic_{\{g>1\}}$ shows $0\le g\le1$ $\rho$-a.e. Rewriting,
$\int f(1-g)\,d\nu=\int fg\,d\mu$. Take $f=\indic_{\{g=1\}}$: $0=\mu(g=1)$, hence
$\nu(g=1)=0$ since $\nu\ll\mu$. On $\{g<1\}$ put $h=g/(1-g)$; for $f=\indic_A/(1-g)$
(truncated and passed to the limit by the MCT) we obtain $\nu(A)=\int_Ah\,d\mu$.
Uniqueness: if $\int_Ah_1d\mu=\int_Ah_2d\mu$ for all $A$, take $A=\{h_1>h_2\}$ and
$\{h_1<h_2\}$ to see both have $\mu$-measure $0$.
\end{proof}

When $\nu=\mu_X$ is the law of $X$ and $\mu$ is Lebesgue measure, $h$ is exactly the density
of Chapter~\ref{ch:cont}; when $\mu$ is counting measure on a countable set, $h$ is the PMF.

\begin{theorem}[Lebesgue decomposition; structure of distributions]
Every distribution function $F$ on $\R$ decomposes as
$F=\alpha_dF_d+\alpha_{ac}F_{ac}+\alpha_sF_s$ with non-negative weights summing to $1$, where
$F_d$ is the CDF of a discrete law (atoms), $F_{ac}$ has a density, and $F_s$ is continuous
but singular (concentrated on a Lebesgue null set, like the Cantor distribution). (Cited.)
\end{theorem}

This is the complete classification promised in Chapter~\ref{ch:cont}: discrete,
absolutely continuous, and mixtures, together with the exotic singular continuous part.

\section{Characteristic functions: inversion, uniqueness, continuity}

Recall $\varphi_X(t)=\Exp[e^{itX}]$. We prove the cited statements of
Theorem~\ref{thm:cfprops}.

\begin{theorem}[Derivatives and moments]\label{thm:cfmoments}
If $\Exp\abs X^k<\infty$ then $\varphi_X$ is $k$ times continuously differentiable and
$\varphi_X^{(k)}(0)=i^k\Exp[X^k]$. If $\Exp[X^2]<\infty$ then
$\varphi_X(t)=1+it\Exp[X]-\tfrac12t^2\Exp[X^2]+o(t^2)$.
\end{theorem}
\begin{proof}
Differentiating $e^{itX}$ $k$ times in $t$ gives $(iX)^ke^{itX}$, bounded in modulus by $\abs X^k$,
which is integrable. By dominated convergence (differences quotients are bounded by
$\abs X^k$ via the mean value theorem), differentiation under the expectation is permitted,
giving $\varphi^{(k)}(t)=\Exp[(iX)^ke^{itX}]$, continuous in $t$ by dominated convergence
again. Set $t=0$. The expansion is Taylor's theorem for $\varphi$ at $0$ with Peano remainder.
\end{proof}

\begin{theorem}[L\'evy inversion formula; uniqueness]\label{thm:inversion}
Let $F$ be the CDF of $X$ and $\varphi$ its characteristic function. For $a<b$ continuity
points of $F$,
\[
F(b)-F(a)=\lim_{T\to\infty}\frac1{2\pi}\int_{-T}^T\frac{e^{-ita}-e^{-itb}}{it}\,\varphi(t)\,dt .
\]
Consequently $\varphi$ determines the distribution of $X$. If $\int\abs{\varphi(t)}dt<\infty$,
then $X$ has a bounded continuous density $f(x)=\frac1{2\pi}\int e^{-itx}\varphi(t)\,dt$.
\end{theorem}
\begin{proof}
The kernel $\frac{e^{-ita}-e^{-itb}}{it}=\int_a^be^{-itx}dx$ has modulus at most $b-a$. By Fubini
(bounded integrand on $[-T,T]\times\Omega$),
\[
I_T:=\int_{-T}^T\frac{e^{-ita}-e^{-itb}}{it}\varphi(t)\,dt=\Exp\Bigl[\int_{-T}^T\frac{e^{it(X-a)}-e^{it(X-b)}}{it}\,dt\Bigr].
\]
For real $c$, $\int_{-T}^T\frac{e^{itc}}{it}dt=\int_{-T}^T\frac{\sin tc}t\,dt-i\int_{-T}^T\frac{\cos tc}t\,dt
=2\,\mathrm{Si}(Tc)$ (the second integral vanishes by oddness, in the principal-value
sense; the combination of the two terms is integrable at $0$), where
$\mathrm{Si}(y)=\int_0^y\frac{\sin u}u\,du$ is odd, bounded, and
$\mathrm{Si}(y)\to\pm\pi/2$ as $y\to\pm\infty$ (the Dirichlet integral, cited from
calculus). Hence the inside of the expectation is
$2[\mathrm{Si}(T(X-a))-\mathrm{Si}(T(X-b))]$, which is bounded uniformly in $T$ and converges
as $T\to\infty$ to $\pi[\mathrm{sgn}(X-a)-\mathrm{sgn}(X-b)]$. This equals $2\pi$ for
$a<X<b$, $0$ for $X<a$ or $X>b$, and $\pi$ when $X\in\{a,b\}$, an event of
probability $0$ since $a,b$ are continuity points. By bounded convergence,
$I_T\to2\pi\Prob(a<X<b)=2\pi(F(b)-F(a))$.

Uniqueness: $F$ is determined at its continuity points (a dense set), and a right-continuous
function is determined by its values on a dense set. For the density: when $\varphi$ is
integrable the limit is an absolutely convergent integral, so by Fubini
$F(b)-F(a)=\frac1{2\pi}\int\varphi(t)\int_a^be^{-itx}dx\,dt=\int_a^bf(x)\,dx$ with
$f(x)=\frac1{2\pi}\int e^{-itx}\varphi(t)dt$.
\end{proof}

This completes the derivation of the Cauchy characteristic function in
Example~\ref{ex:cauchycf} and the uniqueness theorem for characteristic functions (and, via
the same transform argument for Laplace transforms, for MGFs, which are cited in
Theorem~\ref{thm:mgfprops}).

\begin{theorem}[Bounded continuous functions characterise convergence in distribution]\label{thm:portmanteau}
$X_n\xrightarrow{d}X$ iff $\Exp[g(X_n)]\to\Exp[g(X)]$ for every bounded continuous
$g:\R\to\R$.
\end{theorem}
\begin{proof}
($\Leftarrow$) For continuity point $x$ of $F_X$ and $\epsilon>0$, let $g_\epsilon$ be $1$ on
$(-\infty,x-\epsilon]$, $0$ on $[x,\infty)$, linear between, and $h_\epsilon$ the analogous
function between $x$ and $x+\epsilon$. Then $g_\epsilon\le\indic_{(-\infty,x]}\le h_\epsilon$,
so $\Exp g_\epsilon(X_n)\le F_n(x)\le\Exp h_\epsilon(X_n)$; letting $n\to\infty$ and then
$\epsilon\downarrow0$ gives $F_n(x)\to F(x)$.

($\Rightarrow$) Let $g$ be bounded by $B$ and continuous, and $\epsilon>0$. Choose $K$ with
$\pm K$ continuity points of $F_X$ and $\Prob(\abs X>K)<\epsilon$. Since $g$ is uniformly
continuous on $[-K,K]$, there is a partition $-K=x_0<\dots<x_m=K$ into continuity points of
$F_X$ such that $g$ varies by less than $\epsilon$ on each piece; let
$s=\sum_jg(x_j)\indic_{(x_{j-1},x_j]}$. Then $\abs{\Exp g(X_n)-\Exp s(X_n)}\le\epsilon+B\Prob(\abs{X_n}>K)$,
and $\Exp s(X_n)=\sum_jg(x_j)(F_n(x_j)-F_n(x_{j-1}))\to\Exp s(X)$. Also
$\Prob(\abs{X_n}>K)\to\Prob(\abs X>K)<\epsilon$. Combining the estimates for $X_n$ and $X$ gives
$\limsup\abs{\Exp g(X_n)-\Exp g(X)}\le C\epsilon$ for a constant $C$ depending only on $B$.
\end{proof}

\begin{theorem}[L\'evy's continuity theorem]\label{thm:levy}
\begin{enumerate}[label=(\roman*)]
\item If $X_n\xrightarrow{d}X$ then $\varphi_{X_n}(t)\to\varphi_X(t)$ for every $t$.
\item If $\varphi_{X_n}(t)\to\varphi(t)$ for every $t$ and $\varphi$ is continuous at $0$, then
$\varphi$ is the characteristic function of some random variable $X$ and $X_n\xrightarrow{d}X$.
\end{enumerate}
\end{theorem}
\begin{proof}
(i) $x\mapsto\cos tx$ and $\sin tx$ are bounded continuous; apply
Theorem~\ref{thm:portmanteau}.

(ii) \emph{Tightness.} For $u>0$, by Fubini,
\[
\frac1u\int_{-u}^u\bigl(1-\varphi_n(t)\bigr)dt=\Exp\Bigl[\frac1u\int_{-u}^u(1-e^{itX_n})dt\Bigr]
=\Exp\Bigl[2\Bigl(1-\frac{\sin uX_n}{uX_n}\Bigr)\Bigr]\ge\Prob\bigl(\abs{X_n}\ge2/u\bigr),
\]
because the integrand is $\ge0$ always and $\ge2(1-\tfrac1{\abs{uX_n}})\ge1$ when $\abs{uX_n}\ge2$.
Given $\epsilon>0$, continuity of $\varphi$ at $0$ with $\varphi(0)=1$ gives $u$ such that
$\frac1u\int_{-u}^u\abs{1-\varphi(t)}dt<\epsilon$; since $\varphi_n\to\varphi$ pointwise and
boundedly, $\frac1u\int_{-u}^u(1-\varphi_n)\to\frac1u\int_{-u}^u(1-\varphi)$. So for large $n$,
$\Prob(\abs{X_n}\ge2/u)<2\epsilon$: the family is \emph{tight} (no mass escapes to infinity).

\emph{Subsequences.} By Helly's selection theorem (cited), every subsequence of $(F_n)$
has a further subsequence converging at all continuity points to a function $G$ that is
non-decreasing and right-continuous, and tightness ensures $G$ is a genuine CDF (limits
$0$ and $1$ at $\mp\infty$). By (i) the characteristic function of $G$ is the pointwise
limit $\varphi$. By uniqueness (Theorem~\ref{thm:inversion}), $G$ is the same CDF for every
subsequence. A sequence all of whose subsequences have a further subsequence converging to
the same limit must itself converge to that limit, so $F_n\to G$ at its continuity points.
\end{proof}

With this, the central limit theorem proof of Chapter~\ref{ch:limits} is complete.

\section{Conditional expectation, rigorously}

Let $Y$ be an integrable random variable on $(\Omega,\F,\Prob)$ and $\mathcal G\subseteq\F$ a
sub-$\sigma$-algebra (``the information we are allowed to use'').

\begin{definition}[Conditional expectation]
A random variable $Z=\Exp[Y\given\mathcal G]$ is a \emph{conditional expectation of $Y$ given
$\mathcal G$} if (i) $Z$ is $\mathcal G$-measurable and integrable, and (ii)
$\Exp[Y\indic_G]=\Exp[Z\indic_G]$ for every $G\in\mathcal G$. For a random variable $W$ we write
$\Exp[Y\given W]=\Exp[Y\given\sigma(W)]$.
\end{definition}

\begin{theorem}[Existence and uniqueness]\label{thm:condexpexists}
A conditional expectation exists and is unique up to almost sure equality.
\end{theorem}
\begin{proof}
\emph{Uniqueness.} If $Z,Z'$ both satisfy (i)--(ii), then $A=\{Z>Z'\}\in\mathcal G$ and
$\Exp[(Z-Z')\indic_A]=0$, forcing $\Prob(A)=0$ (a strictly positive function on $A$ with zero
integral); similarly for $\{Z<Z'\}$.

\emph{Existence.} For $Y\ge0$ define the finite measure $\nu(G)=\Exp[Y\indic_G]$ on
$(\Omega,\mathcal G)$. Then $\nu\ll\Prob|_{\mathcal G}$ (if $\Prob(G)=0$ then $\nu(G)=0$). By
Radon--Nikodym (Theorem~\ref{thm:rn}) there is a $\mathcal G$-measurable $Z=d\nu/d\Prob|_{\mathcal G}\ge0$
with $\nu(G)=\Exp[Z\indic_G]$, which is property (ii). For general $Y$ apply this to
$Y^+$ and $Y^-$ and subtract.
\end{proof}

\begin{theorem}[Properties]\label{thm:condexpprops}
Let $Y,Y'$ be integrable and $\mathcal H\subseteq\mathcal G\subseteq\F$.
\begin{enumerate}[label=(\roman*)]
\item (Linearity, monotonicity) $\Exp[aY+bY'\given\mathcal G]=a\Exp[Y\given\mathcal G]+b\Exp[Y'\given\mathcal G]$; $Y\le Y'$ implies $\Exp[Y\given\mathcal G]\le\Exp[Y'\given\mathcal G]$ a.s.
\item (Tower) $\Exp[\Exp[Y\given\mathcal G]\given\mathcal H]=\Exp[Y\given\mathcal H]$; in particular $\Exp[\Exp[Y\given\mathcal G]]=\Exp[Y]$.
\item (Taking out what is known) If $W$ is $\mathcal G$-measurable and $WY$ is integrable, then $\Exp[WY\given\mathcal G]=W\Exp[Y\given\mathcal G]$.
\item (Independence) If $Y$ is independent of $\mathcal G$ then $\Exp[Y\given\mathcal G]=\Exp[Y]$.
\item (Jensen) If $\varphi$ is convex and $\varphi(Y)$ integrable, $\varphi(\Exp[Y\given\mathcal G])\le\Exp[\varphi(Y)\given\mathcal G]$.
\end{enumerate}
\end{theorem}
\begin{proof}
(i) The right-hand sides are $\mathcal G$-measurable and satisfy (ii) of the definition by linearity
of expectation; uniqueness identifies them. Monotonicity: apply the defining property on
$G=\{\Exp[Y\given\mathcal G]>\Exp[Y'\given\mathcal G]\}$. (ii) For $H\in\mathcal H\subseteq\mathcal G$,
$\Exp[\Exp[Y\given\mathcal G]\indic_H]=\Exp[Y\indic_H]$; so $\Exp[Y\given\mathcal H]$ also satisfies the
defining property for $\Exp[Y\given\mathcal G]$ restricted to $\mathcal H$. (iii) For
$W=\indic_{G_0}$, $\Exp[\indic_{G_0}Z\indic_G]=\Exp[Z\indic_{G_0\cap G}]=\Exp[Y\indic_{G_0\cap G}]$;
extend to simple $W$ by linearity, to $W\ge0$ by MCT, and to general $W$ by splitting and
dominated convergence. (iv) The constant $\Exp[Y]$ is $\mathcal G$-measurable and
$\Exp[Y\indic_G]=\Exp[Y]\Prob(G)$ by independence. (v) A convex function is the supremum of
countably many affine functions $\ell_j(y)=a_jy+b_j$; $\Exp[\varphi(Y)\given\mathcal G]\ge a_j\Exp[Y\given\mathcal G]+b_j$
a.s.\ for each $j$, and take the supremum over the countable family.
\end{proof}

\begin{theorem}[Orthogonal projection]\label{thm:projection}
If $\Exp[Y^2]<\infty$ then $\Exp[Y\given\mathcal G]$ is the unique (a.s.) minimiser of
$\Exp[(Y-Z)^2]$ over $\mathcal G$-measurable $Z$ with $\Exp[Z^2]<\infty$, and
$Y-\Exp[Y\given\mathcal G]$ is orthogonal to all such $Z$: $\Exp[(Y-\Exp[Y\given\mathcal G])Z]=0$.
\end{theorem}
\begin{proof}
For bounded $\mathcal G$-measurable $Z$, orthogonality is the defining property with $\indic_G$
replaced by $Z$ (linearity and MCT as in (iii) above). For square-integrable $Z$ approximate
by bounded truncations $Z_n=Z\indic_{\{\abs Z\le n\}}$: $Z_n\to Z$ in $L^2$, and
$\abs{\Exp[(Y-\hat Y)(Z-Z_n)]}\le\norm{Y-\hat Y}_2\norm{Z-Z_n}_2\to0$ by Cauchy--Schwarz,
where $\hat Y=\Exp[Y\given\mathcal G]\in L^2$ by conditional Jensen. The Pythagorean identity
$\Exp[(Y-Z)^2]=\Exp[(Y-\hat Y)^2]+\Exp[(\hat Y-Z)^2]$ then gives the minimising property. This
completes the ``limiting argument'' invoked in the proof of Theorem~\ref{thm:bestpred}.
\end{proof}

\begin{nuance}{conditioning on a null event is conditioning on a $\sigma$-algebra}
Every conditional expectation is a random variable measurable with respect to a
$\sigma$-algebra; ``$\Exp[Y\given W=w]$'' is the function $g$ with
$\Exp[Y\given W]=g(W)$ (Doob--Dynkin), defined only up to $\mu_W$-null sets of $w$. The
Borel--Kolmogorov paradox of Chapter~\ref{ch:condexp} arises because ``conditioning on the great
circle'' means different $\sigma$-algebras (generated by longitude versus by latitude).
When one needs the whole conditional \emph{law} of $Y$ given $W=w$, jointly measurable in
$w$, one uses a \emph{regular conditional distribution}, which exists for real-valued $Y$
(cited). For discrete $W$, or jointly continuous $(Y,W)$ with $f_W>0$, it is exactly the
elementary construction of Chapter~\ref{ch:condexp}.
\end{nuance}

\section{Summary of the foundations}

A probability space $(\Omega,\F,\Prob)$ is a measure space of total mass $1$. A random
variable is a measurable function, and its law is the push-forward measure on $\R$,
determined by the CDF. Expectation is the Lebesgue integral. The MCT, Fatou and DCT justify
exchanging limits and expectations, and Tonelli/Fubini justify exchanging integrals. Independence
is the product structure of the joint law. The zero--one law shows tail events are
deterministic. Conditional expectation is a $\sigma$-algebra-measurable projection.
Characteristic functions encode laws uniquely and convergence in distribution as
pointwise convergence of the transforms. Everything proved in earlier chapters for discrete
and density-based random variables holds in this generality.

\section*{Exercises}
\addcontentsline{toc}{section}{Exercises}

\begin{enumerate}
\item Show that $\sigma(X)$ is a $\sigma$-algebra and that $X$ is $\mathcal G$-measurable iff
$\sigma(X)\subseteq\mathcal G$.
\item Show that the Borel $\sigma$-algebra on $\R$ is generated by (a) the intervals $(a,b]$,
(b) the half-lines $(-\infty,x]$ with $x$ rational.
\item Prove that a finitely additive $\mu$ on a $\sigma$-algebra with $\mu(\Omega)<\infty$ is
countably additive iff it is continuous from above at $\emptyset$.
\item Let $X_n\ge0$ be integrable with $X_n\to0$ a.s.\ and $\Exp[X_n]=1$. Show that DCT fails
and find where the domination condition breaks (Example~\ref{ex:conversefail}(ii)).
\item Prove that $\sum_{n=1}^\infty\Exp\abs{X_n}<\infty$ implies $\sum X_n$ converges almost
surely and $\Exp[\sum X_n]=\sum\Exp[X_n]$ (apply the MCT and Fubini to counting measure).
\item Show that if $X,Y$ are independent with $\Exp\abs X<\infty$ then
$\Exp[X\given Y]=\Exp[X]$ using the rigorous definition.
\item Compute the characteristic function of $U(-1,1)$ and verify the inversion
formula at $x=0$.
\item Prove Fatou's lemma implies the first Borel--Cantelli direction: if
$\sum\Prob(A_n)<\infty$ then $\Prob(A_n\text{ i.o.})=0$. (Apply the MCT to
$\sum_n\indic_{A_n}$.)
\item (Challenge) Prove Helly's selection theorem: every sequence of CDFs has a
subsequence converging at all continuity points of a non-decreasing right-continuous limit
(a diagonal argument along the rationals).
\end{enumerate}
